% Generated draft: requires mathematical and source review.
\documentclass[11pt]{article}
\usepackage[margin=25mm]{geometry}
\usepackage{fontspec}
\setmainfont{Latin Modern Roman}
\usepackage{amsmath,amssymb,mathtools,unicode-math}
\setmathfont{Latin Modern Math}
\usepackage{enumitem,needspace,xurl,hyperref}
\setlength{\parindent}{0pt}
\setlength{\parskip}{5pt}
\setlength{\emergencystretch}{3em}
\newcommand{\N}{\mathbb{N}}
\newcommand{\Z}{\mathbb{Z}}
\newcommand{\Q}{\mathbb{Q}}
\newcommand{\R}{\mathbb{R}}
\newcommand{\C}{\mathbb{C}}
\newcommand{\F}{\mathbb{F}}
\newcommand{\A}{\mathbb{A}}
\newcommand{\PP}{\mathbb{P}}
\newcommand{\E}{\mathbb{E}}
\newcommand{\eps}{\varepsilon}
\newcommand{\dd}{\,\mathrm{d}}
\begin{document}
\section*{UnsolvedMath ID: 30005523}
\textbf{Source problem code:} catalogue entry
\subsection{Definitions and mathematical statement}
For each viscosity parameter \(\varepsilon>0\), let \(u_\varepsilon\) denote the solution asserted in the supplied record to be the unique weak solution of
\[
-\varepsilon\Delta u_\varepsilon+Hu_\varepsilon+\alpha u_\varepsilon
=\langle f,\xi^m\rangle.
\]
Here \(\Delta\), \(H\), \(\alpha\), \(f\), \(\xi\), \(m\), and the pairing \(\langle f,\xi^m\rangle\) retain their source-dependent meanings; the supplied record does not define them. In particular, it does not specify the underlying domain, the operator domains, the admissible data, the boundary or initial conditions, or the function spaces and test functions defining a weak solution. Existence and uniqueness are premises reported by the source record, not assertions established by this draft.

The extracted question is whether, with the underlying problem data fixed, the full family \(u_\varepsilon\) has a limit
\[
u_\varepsilon\longrightarrow u\qquad\text{as }\varepsilon\downarrow0
\]
in a mathematically appropriate topology, to be identified, and, if so, whether that limit solves the corresponding original transport problem. Formally deleting the viscosity term gives the candidate equation
\[
Hu+\alpha u=\langle f,\xi^m\rangle.
\]
The intended weak formulation and the conditions required of this limit are not supplied. Moreover, the original excerpt refers to the original transport equations in the plural, without identifying that system. Consequently, the displayed candidate equation records only the formal target suggested by the supplied regularized equation; a complete quantified convergence statement and an identification with the entire original transport problem require the missing source context.
\subsection{Short English statement}
Does the asserted family of unique weak solutions converge as the positive viscosity tends to zero, in an appropriate topology, and does its limit solve the original transport problem? The supplied record leaves the domain, operators, data assumptions, solution spaces, and boundary or initial conditions unspecified, so the target remains conditional on recovering that context.
\subsection{Sources}
\begin{itemize}
\item Supplied local registry record: permanent UnsolvedMath ID \(30005523\), title \emph{Vanishing-Viscosity Limits for Transport Equations}, source problem code \texttt{catalogue entry}, in the Oberwolfach Reports Open Problems collection.
\item Supplied report citation: \emph{Tomographic Inverse Problems: Mathematical Challenges and Novel Applications} (2023), Oberwolfach Reports, pp.\ 1105--1194, DOI \texttt{10.4171/owr/2023/21}. Supplied source and literature-source link, labelled ``Oberwolfach problem report'': \url{https://doi.org/10.4171/owr/2023/21}. No individual proposer or precise location of the extracted passage is supplied.
\item Literature assessment dated exactly \texttt{2026-08-22}: the supplied curation labels the problem open and reports that existence and topology of the singular limit, and its passage to the original transport equations, remain unresolved in the stated generality. This is supplied dated curation, not independently verified here; the background expressly requests independent verification of time-sensitive status and cited claims.
\item The registry labels the statement \texttt{corrected\_verified} and reports that the official report was checked directly and that definitions were restored. These are supplied provenance claims, not verification performed for this draft; the complete record provided here still omits the definitions needed for a precise standalone formulation.
\end{itemize}
\end{document}
